\section{Limitations and threats to validity}
\label{sec:limitations}

\paragraph{One GPU, one vendor, one toolkit, one part class.} Every number here
comes from a single RTX 5070 running CUDA Toolkit 13.3 under WSL2. It is a
consumer part, not a datacenter one: 12 GB of GDDR7 rather than HBM, a different
memory system, a different power and clock envelope, and no NVLink. Nothing in this
report has been run on a second card, a second toolkit version, or a second vendor.
Where a claim depends on the hardware, and most of them do, it is a claim about
this configuration.

\paragraph{The negative result is a single GPU finding.} Section \ref{sec:negative}
shows shared memory tiling losing to naive, and names the mechanism: the last level
cache absorbs the redundant traffic that tiling exists to remove. That mechanism is
a statement about the ratio of cache capacity to working set, which means the sign
of the result should flip on a part with a smaller cache or a larger working set,
and I have not shown that. The study that would establish it is a crossover
experiment across three generations, predicting the crossing point from L2 capacity
and confirming it on Ampere, Ada and Blackwell. It needs hardware this machine does
not have, and it is scoped to a future release rather than claimed here. Until
somebody runs it, the correct reading of Section \ref{sec:negative} is "on this
part", not "in general".

\paragraph{Provenance.} Corrections A1 and A4 are closed. Every row of the
committed summary carries a commit hash that resolves and a locked clock, and a
style gate fails the build on any recorded hash that does not resolve. Seven dead
v1 hashes are still allowed through by
\texttt{experiments/results/legacy\_hashes.txt}, and that is deliberate and scoped:
they stamp archived material that is kept in the tree as history, the superseded
rows in \texttt{sweep.jsonl} and the round 01 to 09 profiler pages. No figure in
this report is traceable to one of them.

\paragraph{Contaminated denominators, and what survived them.} Correction A2
recorded setup work inside two vendor timed regions. The wrappers were fixed and
the sweep re-run, so the percentages here are measured against a cuSPARSE plan
built once and reused and against a TRSM baseline whose restore copy is outside the
window. The v1 percentages are not corrected, they are withdrawn: the setup cannot
be subtracted after the fact, and the generator still prints the withdrawal for any
row that names no matrix, so a v1 row can never be read as though it had been
measured this way.

\paragraph{The traffic column is not in the summary yet.} Round 15 to 17 measured
\texttt{dram\_\_bytes.sum} for the SpMV, FFT, scan and reduction rungs, and those
figures are quoted in this report and in the diagnostic log. The
\texttt{dram\_bytes\_sum} column of \texttt{summary.csv} still reads
\texttt{pending ncu round}, so a reader who wants the traffic behind a specific row
has to go to the profiler page rather than to the summary. Re-summarizing the sweep
is what closes that, and it has not been done.

\paragraph{The roofline chart is measured at a different clock from the ladder.}
The ceilings come from a device probe at 2842 MHz; every ladder row is at a 2497
MHz lock. Percent of roof is quoted against the locked clock roof throughout, but
the chart itself is drawn against the probe's ceilings, so a point read off the
figure by eye is not the same quantity as a percent of roof quoted in the text.

\paragraph{What the CI bounds do and do not cover.} The performance regression job
compares a quick sweep against a committed baseline, which catches a large step
change on this machine. It does not establish that a small difference is real: a
quick sweep is fewer repetitions than the full protocol, and it runs on whatever
thermal state the machine happens to be in. The full protocol's five independent
process repeats and their bootstrap intervals are recorded per row and drawn as
error bars on the ladder figure, so a difference between two rows can now be judged
rather than guessed at. The regression job uses both halves of that: a row counts
as a regression only when it is more than the tolerance slower \emph{and} the two
intervals do not overlap.

\paragraph{Claims I am carrying from elsewhere.} The observation that cuBLAS
dispatches Ampere generation kernels for FP16 GEMM on consumer Blackwell is
somebody else's measurement~\cite{cloudrift}, not mine. I use it as context for why
a gap exists at all. I have not reproduced it, and no result in this report depends
on it being true.

\paragraph{Benchmark scope.} The sweep covers four square shapes for GEMM and a
handful of sizes elsewhere. Square shapes are the friendliest case for a fixed tile
and the least interesting one for a dispatch heuristic. Non square, batched, and
strided cases are supported by the API and tested for correctness, but they are not
swept, so nothing here should be read as a performance claim about them.
